\subsection{\texorpdfstring{EXAMPLE I: COMPACT ALGEBRAS OF TYPE \(A_{n}\)}{EXAMPLE I: COMPACT ALGEBRAS OF TYPE A\_\{n\}}}

Let V be a finite dimensional complex vector space and \(\Phi\) a separating positive hermitian form on V. The unitary group associated to \(\Phi\) (cf.~Algebra, Chap. IX) is the subgroup \(\mathbf{U}(\Phi)\) of \(\mathbf{GL}(V)\) consisting of the automorphisms of the complex Hilbert space \((V,\Phi)\) ; this is a (real) Lie subgroup of the group \(\mathbf{GL}(V)\) , whose Lie algebra is the subalgebra \(\mathfrak{u}(\Phi)\) of the real Lie algebra \(\mathfrak{gl}(V)\) consisting of the endomorphisms x of V such that \(x^{*} = -x\) (Chap. III, §3, no. 10, Cor. 2 of Prop. 37), where \(x^{*}\) denotes the adjoint of x relative to \(\Phi\) . Since the group \(\mathbf{U}(\Phi)\) is compact (§1, no. 1), \(\mathfrak{u}(\Phi)\) is a compact real Lie algebra. Similarly, the special unitary group \(\mathbf{SU}(\Phi) = \mathbf{U}(\Phi) \cap \mathbf{SL}(V)\) is a compact Lie subgroup of \(\mathbf{SL}(V)\) , whose Lie algebra is \(\mathfrak{su}(\Phi) = \mathfrak{u}(\Phi) \cap \mathfrak{sl}(V)\) .

When \(V = C^{n}\) and \(\varPhi\) is the usual hermitian form (for which the canonical basis of \(C^{n}\) is orthonormal), we write \(\mathbf{U}(n, \mathbf{C})\) , \(\mathbf{SU}(n, \mathbf{C})\) , \(\mathfrak{u}(n, \mathbf{C})\) , \(\mathfrak{s}\mathfrak{u}(n, \mathbf{C})\) instead of \(\mathbf{U}(\varPhi)\) , \(\mathbf{SU}(\varPhi)\) , \(\mathfrak{u}(\varPhi)\) , \(\mathfrak{s}\mathfrak{u}(\varPhi)\) . The elements of \(\mathbf{U}(n, \mathbf{C})\) (resp. \(\mathfrak{u}(n, \mathbf{C})\) )

are the matrices \(A \in \mathrm{M}_{n}(\mathbf{C})\) such that \(A.^{t}\bar{A} = I_{n}\) (resp. \(A = -^{t}\bar{A}\) ), which are said to be unitary (resp. anti-hermitian).

PROPOSITION 4. a) The compact real forms of the complex Lie algebra \(\mathfrak{sl}(V)\) are the algebras \(\mathfrak{su}(\Phi)\) , where \(\Phi\) belongs to the set of separating positive hermitian forms on the complex vector space \(V\) .

\begin{enumerate}
\def\labelenumi{\alph{enumi})}
\setcounter{enumi}{1}
\tightlist
\item
  The algebras \(\mathfrak{u}(\varPhi)\) are the compact real forms of \(\mathfrak{gl}(\mathrm{V})\) .
\end{enumerate}

Let \(\varPhi\) be a separating positive hermitian form on V. For all \(x\in\mathfrak{gl}(V)\) , put \(\sigma(x)=-x^{*}\) (where \(x^{*}\) is the adjoint of \(x\) relative to \(\varPhi\) ). Then \(\sigma\) satisfies conditions (2) of Prop. 1 of no. 1, so the set \(\mathfrak{u}(\varPhi)\) (resp. \(\mathfrak{su}(\varPhi)\) ) of fixed points of \(\sigma\) on \(\mathfrak{gl}(V)\) (resp. \(\mathfrak{sl}(V)\) ) is a compact real form of \(\mathfrak{gl}(V)\) (resp. \(\mathfrak{sl}(V)\) ). Since \(\mathbf{GL}(V)\) operates transitively on the set of separating positive hermitian forms on V (Algebra, Chap. IX) and on the set of compact real forms of \(\mathfrak{sl}(V)\) (no. 3, Th. 1 and Chap. VIII, §13, no. 1 (VII)), Prop. 4 is proved.

COROLLARY. Every compact simple real Lie algebra of type \(\mathrm{A}_n\) ( \(n \geq 1\) ) is isomorphic to \(\mathfrak{su}(n + 1, \mathbf{C})\) .

Indeed, every complex Lie algebra of type \(\mathbf{A}_n\) is isomorphic to \(\mathfrak{sl}(n + 1,\mathbf{C})\) (Chap. VIII, §13, no. 1).

Remarks. 1) We have \(\mathfrak{gl}(\mathrm{V}) = \mathfrak{sl}(\mathrm{V}) \times \mathbf{C}.1_{\mathrm{V}}\) , \(\mathfrak{u}(\varPhi) = \mathfrak{su}(\varPhi) \times \mathbf{R}.i1_{\mathrm{V}}\) ; the compact real forms of \(\mathfrak{gl}(\mathrm{V})\) are the \(\mathfrak{su}(\varPhi) \times \mathbf{R}. \alpha 1_{\mathrm{V}}, \alpha \in \mathbf{C}^*\) .

\begin{enumerate}
\def\labelenumi{\arabic{enumi})}
\setcounter{enumi}{1}
\tightlist
\item
  If the complex Lie algebra \(\mathfrak{a} = \mathfrak{sl}(n,\mathbf{C})\) is equipped with the splitting and Chevalley system introduced in Chap. VIII, §13, no. 1 (IX), then, with the notations in no. 2,
\end{enumerate}

\[
\mathfrak {a} _ {u} = \mathfrak {s u} (n, \mathbf {C}), \quad \mathfrak {a} _ {0} = \mathfrak {s l} (n, \mathbf {R}), \quad \mathfrak {a} _ {u} \cap \mathfrak {a} _ {0} = \mathfrak {o} (n, \mathbf {R}).
\]

